\documentclass[12pt,a4paper]{article}
\usepackage[utf8]{inputenc}
\usepackage[spanish]{babel}
\usepackage{booktabs}
\usepackage{xcolor}
\usepackage{geometry}
\usepackage{enumitem}
\usepackage{parskip}
\usepackage{float}
\usepackage{graphicx}
\usepackage{amsmath}
\usepackage{longtable}
\newcommand{\textbffeat}{\textbf{feat}}

\geometry{margin=2.5cm}

\title{Línea Temporal de Errores -- Repositorio Tesis}
\author{Fabián}
\date{\today}

\begin{document}

\maketitle

\section*{Errores en collision\_erosion\_bed.input (LAMMPS)}

\subsection*{2026-03-22 -- Problemas con halt y wave detection}

\begin{table}[H]
	\centering
	\begin{tabular}{@{}llp{9cm}@{}}
		\toprule
		\textbf{Hora} & \textbf{Commit} & \textbf{Error y Solución}                                                                                                                                    \\
		\midrule
		13:47         & e3aa79d         & \textbf{Error en step y defaults proc:} Uso incorrecto de variables en collision\_erosion\_bed. Corrección de defaults de proc y manejo de step.             \\
		13:55         & 366cb55         & \textbf{Sintaxis booleana incorrecta:} El if en el loop de trampa usaba lógica booleana incorrecta. Se corrigió la sintaxis.                                 \\
		04:30         & 1be9177         & \textbf{Sin loop de trampa:} El halt de onda cortaba la simulación pero no había forma limpia de salir. Se agregó loop de trampa + salida limpia.            \\
		04:39         & 7afdaef         & \textbf{Could not find compute ID:} El primer uncompute fallaba porque check\_wave no existía aún. Se crea dummy antes del label trampa.                     \\
		04:49         & 7f73005         & \textbf{Could not find compute ID + Compute not current:} El approach de uncompute/recompute fallaba entre runs. Solución KISS: usar file\_exists flag.      \\
		05:06         & 7caef10         & \textbf{Misma solución KISS:} Se refina el approach de flag por archivo para salir del loop tras halt.                                                       \\
		05:09         & 7caef10         & \textbf{Permisos en /tmp:} El cluster no permite escribir en /tmp. Se cambia a /dev/shm para el flag.                                                        \\
		05:22         & 5439e8c         & \textbf{Halt en todas las simulaciones:} El halt se activaba para todas, incluyendo las no críticas. Solo debe afectar vel=-50 y move$\neq$0.                \\
		05:38         & 0bf5145         & \textbf{Compute not current between runs:} El approach con v\_stop\_cond fallaba. Nueva estrategia: comparar step\_before vs step\_after para detectar halt. \\
		\bottomrule
	\end{tabular}
\end{table}

\section*{Errores en CARBOXI (Python/Polars)}

\subsection*{2026-03-25 -- Empaquetamiento y cluster headless}

\begin{table}[H]
	\centering
	\begin{tabular}{@{}llp{9cm}@{}}
		\toprule
		\textbf{Fecha} & \textbf{Commit} & \textbf{Error y Solución}                                                                                                                                                                                                                                                       \\
		\midrule
		2026-03-25     & 272950f         & \textbf{hatchling crea .pth vacío en editable install:} PEP 660 editable install con hatchling creaba un archivo \_carboxi.pth vacío en site-packages. Python no encontraba el módulo. Solución: cambiar build backend a setuptools que maneja editable installs correctamente. \\
		2026-03-25     & 9ede3bc         & \textbf{ovito segfault en cluster headless:} ovito intenta conectarse a un display gráfico al importarse. En cluster sin monitor → segfault. Solución: arrancar Xvfb (display virtual) antes de importar ovito.                                                                 \\
		2026-03-25     & ac6d996         & \textbf{GLIBC conflict con Xvfb:} conda modifica LD\_LIBRARY\_PATH con su propio glibc. Xvfb y tools del sistema fallan con \texttt{undefined symbol: \_\_libc\_pthread\_init}. Solución: guardar LD\_LIBRARY\_PATH antes de activar conda, restaurarlo antes de Xvfb.          \\
		2026-03-25     & ac6d996         & \textbf{Stale Xvfb lock files:} Si Xvfb muere, deja /tmp/.X99-lock que impide arrancar uno nuevo. Solución: \texttt{rm -f /tmp/.X99-lock /tmp/.X11-unix/X99} antes de arrancar.                                                                                                 \\
		2026-03-25     & ac6d996         & \textbf{Editable install falla desde Buffer/:} python -m CARBOXI fallaba con "No module named CARBOXI" desde directorios de simulación. Funcionaba desde /tmp y \~. Solución: ejecutar CARBOXI desde \~{} en el script.                                                         \\
		2026-03-25     & ---             & \textbf{pip vs Python distintos:} pip apuntaba a Python 3.13 del sistema, python era 3.11 del conda env. pip install instalaba en el lugar equivocado. Solución: usar \texttt{python -m pip install} en vez de \texttt{pip install}.                                            \\
		\bottomrule
	\end{tabular}
\end{table}

\subsection*{2026-03-25 -- Método de análisis hardcodeado}

\begin{table}[H]
	\centering
	\begin{tabular}{@{}llp{9cm}@{}}
		\toprule
		\textbf{Fecha} & \textbf{Commit} & \textbf{Error y Solución}                                                                                                                                                                                                                                                                                                               \\
		\midrule
		2026-03-25     & ac6d996         & \textbf{Método hardcodeado a taylor:} El método de análisis de erosión estaba hardcodeado a ``taylor'' en pipeline.py. Los archivos de salida siempre tenían ``\_filtrado\_taylor'' aunque se quería usar ``v1''. Solución: agregar campo \texttt{method} a config.toml [filtros] y usar \texttt{config.filters.method} en pipeline.py. \\
		\bottomrule
	\end{tabular}
\end{table}

\subsection*{2026-03-25/26 -- Filtro de erosión inter-timestep}

\begin{table}[H]
	\centering
	\begin{tabular}{@{}llp{9cm}@{}}
		\toprule
		\textbf{Fecha} & \textbf{Commit} & \textbf{Error y Solución}                                                                                                                                                                                                                                                                                                          \\
		\midrule
		2026-03-25     & 85d4c48         & \textbf{Filtro fz caótico:} El filtro por $\sum fz \approx 0$ no funcionaba bien (valores erráticos). Solución: reemplazar por comparación de $|v_{cm}|$ entre dumps consecutivos (por ID de partícula en común).                                                                                                                  \\
		2026-03-25     & 85d4c48         & \textbf{Clusters de tamaño 1 filtran partículas correctamente eyectadas:} Cuando una partícula está sola en un cluster, $|\Delta v_{cm}| \approx 0$ (se compara consigo misma). El filtro inter-timestep la elimina como ``estacionaria''. Bug pendiente.                                                                          \\
		2026-03-26     & 62590cf         & \textbf{Partículas oscilan alrededor de max\_height:} Partículas a 2--6 radios sobre el umbral pasan brevemente por debajo de $z > max\_height$, generando gaps en el resultado. Solución: filtro por CM z del cluster en vez de z individual. No resolvió el problema porque los clusters de tamaño 1 tienen CM z = z individual. \\
		2026-03-26     & 62590cf         & \textbf{ID 61991 SÍ está en el dump de LAMMPS:} Verificado con \texttt{zgrep} en el cluster. ID 61991 está en output.900000.gz con z=107.8 µm y vz=0.1495 m/s. El gap es causado por CARBOXI que filtra la partícula, no por LAMMPS.                                                                                               \\
		2026-03-26     & 28546c0         & \textbf{Config duplicado:} Dos config.toml (cluster/ y CARBOXI/) con overlap y valores distintos (material, tolerancia). Solución: unificar en un solo config.toml en cluster/ con \texttt{--config} argumento para CARBOXI.                                                                                                       \\
		2026-03-26     & 1c44b0b         & \textbf{Fix gaps definitivo: tracking acumulativo de IDs:} Una vez detectada como erosionada, la partícula se mantiene en todos los frames siguientes. Si oscila por debajo de max\_height, se agrega desde el dump raw. Resultado: 18/18 simulaciones con 0 gaps (antes: 20--85 gaps).                                            \\
		\bottomrule
	\end{tabular}
\end{table}

\subsection*{2026-03-26 -- Comparar filtros}

\begin{table}[H]
	\centering
	\begin{tabular}{@{}llp{9cm}@{}}
		\toprule
		\textbf{Fecha} & \textbf{Commit} & \textbf{Error y Solución}                                                                                                                                                                                                                          \\
		\midrule
		2026-03-26     & 5826f4f         & \textbf{Keyword args incompatibles:} Las funciones de filtro tenían parámetros con prefijo \_ (\_prev\_dump, \_tol). main.py los llamaba sin prefijo. Solución: quitar prefijos de los parámetros.                                                 \\
		2026-03-26     & 5f3abed         & \textbf{Schema mismatch en concat:} \texttt{pl.concat(..., how="diagonal")} fallaba porque \texttt{pl.lit(i)} creaba Int32 inconsistente. Solución: \texttt{pl.lit(i, dtype=pl.Int64)} y \texttt{how="diagonal\_relaxed"}.                         \\
		2026-03-26     & 7fbc8cd         & \textbf{DataFrame vacío de F3C:} Cuando F3C no tiene clusters balísticos, devuelve \texttt{pl.DataFrame()} vacío sin columnas. \texttt{df.filter(pl.col("type"))} falla. Solución: check de \texttt{df.is\_empty()} al inicio de compute\_metrics. \\
		\bottomrule
	\end{tabular}
\end{table}

\subsection*{2026-03-28/30 -- Refactorización y visualización}

\begin{table}[H]
	\centering
	\begin{tabular}{@{}llp{9cm}@{}}
		\toprule
		\textbf{Fecha} & \textbf{Commit} & \textbf{Error y Solución}                                                                                                                                                                                                                                                    \\
		\midrule
		2026-03-30     & 46c10a8         & \textbf{Refactor main.py:} Función \texttt{run\_simulation()} de 210 líneas hacía todo (filtrado, métricas, tracking, gaps, CSV, plots, tablas). Solución: extracción en \texttt{\_process\_config()}, \texttt{\_apply\_tracking()}, \texttt{\_build\_summary()} (SRP, DRY). \\
		2026-03-30     & b33a2d9         & \textbf{Cuarto plot faltante:} Solo había 3 subplots (n\_bed, n\_clusters, gaps acumulados). Solución: agregar \texttt{compute\_gaps\_per\_raw()} y cuarto plot de gaps por step (no acumulados).                                                                            \\
		2026-03-30     & b33a2d9         & \textbf{Ylabel de subplot (d) se superpone con (c):} Espacio insuficiente entre columnas. Solución: agregar \texttt{wspace=0.28} en \texttt{subplots\_adjust()}.                                                                                                             \\
		2026-03-30     & 6c7a324         & \textbf{Partículas congeladas no incluidas:} El filtro inicial solo capturaba partículas con $z > bed\_height$ sin considerar partículas estáticas. Solución: incluir partículas congeladas en el filtro.                                                                    \\
		2026-03-30     & c570419         & \textbf{Revert main.py:} La reescritura completa introdujo bugs. Se restauró versión anterior y se aplicaron cambios incrementales.                                                                                                                                          \\
		2026-03-29     & 974d93c         & \textbf{Estructura de archivos desordenada:} Archivos sueltos en raíz del proyecto. Solución: reorganizar en carpetas lógicas.                                                                                                                                               \\
		\bottomrule
	\end{tabular}
\end{table}

\subsection*{2026-03-15 -- Errores de tipos en Parquet}

\begin{table}[H]
	\centering
	\begin{tabular}{@{}llp{9cm}@{}}
		\toprule
		\textbf{Hora} & \textbf{Commit} & \textbf{Error y Solución}                                                                                               \\
		\midrule
		16:32         & 4967def         & \textbf{Error Object en Parquet:} Los tipos no se forzaban correctamente. Se agregó forzado de tipos en result\_writer. \\
		16:55         & df6bb21         & \textbf{Tipos Object en concat:} Schema no se aplicaba en concat. Se usa schema explícito en concat.                    \\
		17:12         & e202d6e         & \textbf{Tipos Object tras concat:} Conversión de tipos insuficiente. Se convierten tipos antes de concat.               \\
		17:31         & 0b9bdd7         & \textbf{Errores de tipo en cast:} strict=True fallaba. Se usa strict=False en cast.                                     \\
		17:34         & 2201d73         & \textbf{Columnas Object:} Polars convertía tipos a Object. Se reemplazan con null del tipo correcto.                    \\
		\bottomrule
	\end{tabular}
\end{table}

\subsection*{2026-03-15 -- OVITO y plot}

\begin{table}[H]
	\centering
	\begin{tabular}{@{}llp{9cm}@{}}
		\toprule
		\textbf{Hora} & \textbf{Commit} & \textbf{Error y Solución}                                                                                                                      \\
		\midrule
		15:42         & 603b37d         & \textbf{Nombres de propiedades OVITO incorrectos:} 'Cluster ID' no existía. Se cambió a 'Cluster Identifier'. Orden de análisis incorrecto.    \\
		17:39         & 7e18039         & \textbf{plot.py sin proyectil\_erosion:} El nuevo algoritmo de erosión eliminó proyectil\_erosion pero plot.py lo seguía usando. Se actualizó. \\
		\bottomrule
	\end{tabular}
\end{table}

\section*{Errores en Pipeline y Lógica}

\subsection*{2026-03-22 -- Config y scripting}

\begin{table}[H]
	\centering
	\begin{tabular}{@{}llp{9cm}@{}}
		\toprule
		\textbf{Hora} & \textbf{Commit} & \textbf{Error y Solución}                                                                                                                                     \\
		\midrule
		18:33         & 23fbeca         & \textbf{Doble for perdido:} La refactorización a config.toml perdió el producto cartesiano original de moves. Solo quedaba ymove=0. Se restauró el doble for. \\
		\bottomrule
	\end{tabular}
\end{table}

\subsection*{2026-03-25 -- Simplificación CARBOXI}

\begin{table}[H]
	\centering
	\begin{tabular}{@{}llp{9cm}@{}}
		\toprule
		\textbf{Fecha} & \textbf{Commit} & \textbf{Error y Solución}                                                                                                                                                                                                                                                                                           \\
		\midrule
		2026-03-25     & adcab92         & \textbf{Código muerto acumulado (242 líneas):} Método taylor nunca usado, routing de método innecesario, código comentado en pipeline.py y lammps\_reader.py, campos de config sin uso (method, lonja\_radios). Solución: eliminación masiva siguiendo YAGNI. Solo queda el método v1 (z+vz + filtro balístico fz). \\
		\bottomrule
	\end{tabular}
\end{table}

\subsection*{2026-03-20 -- Convert}

\begin{table}[H]
	\centering
	\begin{tabular}{@{}llp{9cm}@{}}
		\toprule
		\textbf{Hora} & \textbf{Commit} & \textbf{Error y Solución}                                                                                              \\
		\midrule
		21:48         & 593caa9         & \textbf{Polars no soporta lista de Paths:} pl.read\_csv(list[Path]) no funciona. Se leyó cada archivo individualmente. \\
		\bottomrule
	\end{tabular}
\end{table}

\subsection*{2026-03-16 -- Detección de erosión}

\begin{table}[H]
	\centering
	\begin{tabular}{@{}llp{9cm}@{}}
		\toprule
		\textbf{Hora} & \textbf{Commit} & \textbf{Error y Solución}                                                                                                        \\
		\midrule
		02:01         & f493b8b         & \textbf{Umlaut de altura sin radios:} La detección de erosión usaba altura sin considerar radios. Se agregaron radios al umbral. \\
		\bottomrule
	\end{tabular}
\end{table}

\subsection*{2026-03-12 -- Criticidad}

\begin{table}[H]
	\centering
	\begin{tabular}{@{}llp{9cm}@{}}
		\toprule
		\textbf{Hora} & \textbf{Commit} & \textbf{Error y Solución}                                                                                                                                  \\
		\midrule
		19:39         & 64e63c8         & \textbf{Criticidad evaluada mal:} Se evaluaba solo la primera carpeta. Todas las simulaciones usaban esos parámetros. Se evaluó por simulación individual. \\
		\bottomrule
	\end{tabular}
\end{table}

\section*{Errores Históricos}

\subsection*{2026-03-07 -- data\_extractor y detector\_ondas}

\begin{table}[H]
	\centering
	\begin{tabular}{@{}llp{9cm}@{}}
		\toprule
		\textbf{Hora} & \textbf{Commit} & \textbf{Error y Solución}                                                                                                                                                                             \\
		\midrule
		20:12         & 3967bc2         & \textbf{Borré parte de buscar\_carpeta():} Error al refactorizar, el código dejó de funcionar. Además, detector\_ondas() tomaba al proyectil como parte de la onda. Se filtró cuando está en el aire. \\
		\bottomrule
	\end{tabular}
\end{table}

\section*{Resumen de Errores por Categoría}

\subsection*{LAMMPS Input (collision\_erosion\_bed.input)}
\begin{itemize}[leftmargin=*]
	\item \textbf{Halt/Wave Detection:} 9 fixes relacionados con halt, compute IDs, permisos, y lógica booleana
	\item \textbf{Main Issues:} Compute not current, Could not find compute ID, permisos en /tmp
\end{itemize}

\subsection*{Python/Polars (CARBOXI)}
\begin{itemize}[leftmargin=*]
	\item \textbf{Tipos Parquet:} 5 fixes para columnas Object y errores de cast
	\item \textbf{OVITO:} Nombres de propiedades y orden de análisis
	\item \textbf{Plot:} Variables eliminadas en refactor
\end{itemize}

\subsection*{Empaquetamiento y Cluster}
\begin{itemize}[leftmargin=*]
	\item \textbf{Build backend:} hatchling crea .pth vacío → setuptools
	\item \textbf{Headless:} ovito segfault sin display → Xvfb
	\item \textbf{GLIBC:} conda corrompe LD\_LIBRARY\_PATH → restaurar antes de Xvfb
	\item \textbf{Lock files:} Xvfb deja locks stale → cleanup manual
	\item \textbf{Editable install:} no funciona desde ciertos directorios → ejecutar desde \~{}
	\item \textbf{pip mismatch:} pip vs python de distintas versiones → usar \texttt{python -m pip}
\end{itemize}

\subsection*{Pipeline}
\begin{itemize}[leftmargin=*]
	\item \textbf{Configuración:} Doble for de moves perdido en refactor
	\item \textbf{Convert:} Polars no soporta lista de Paths
	\item \textbf{Criticidad:} Evaluada solo en primera carpeta
	\item \textbf{Detección:} Umbral sin considerar radios
	\item \textbf{Simplificación:} 242 líneas de código muerto eliminadas (método taylor, routing, código comentado)
	\item \textbf{Filtro fz:} Caótico → reemplazado por comparación CM inter-timestep
	\item \textbf{Gaps:} Partículas oscilan alrededor de max\_height → tracking acumulativo de IDs (18/18 con 0 gaps)
	\item \textbf{Config unificado:} Dos config.toml con overlap → merge en cluster/config.toml
\end{itemize}

\subsection*{Comparar Filtros}
\begin{itemize}[leftmargin=*]
	\item \textbf{Función monolítica:} \texttt{run\_simulation()} de 210 líneas → dividida en 3 funciones (SRP)
	\item \textbf{Estado global:} \texttt{global TOLERANCIA\_V\_CM} → parámetro \texttt{tol} en cadena de llamadas
	\item \textbf{Ovito innecesario:} Se ejecutaba para métricas de filtros individuales → solo en cluster
	\item \textbf{Ylabel superpuesto:} Subplot (d) tapaba con (c) → \texttt{wspace=0.28}
	\item \textbf{Partículas congeladas:} No incluidas en filtro inicial → filtro corregido
\end{itemize}

\subsection*{Data Extractor}
\begin{itemize}[leftmargin=*]
	\item \textbf{Funciones borradas:} buscar\_carpeta() parcialmente eliminada
	\item \textbf{Lógica de onda:} Proyectil incluido incorrectamente en detección
\end{itemize}

\end{document}
